% Historical source (removed by author request): references/_OLD_Übungsaufgaben/Übungsblatt 4 - Funktionen/Blatt4.tex:274
\begin{ExerciseBlock}
\ExerciseHeading{16}{Hyperbolikus-Funktionen}{functions_hyperbolic_functions}
Die Hyperbolikus-Funktionen (Sinus-Hyperbolikus,\newline
Kosinus-Hyperbolikus und Tangens-Hyperbolikus) sind definiert als
\begin{equation*}
	\begin{split}
		\sinh(x) &= \frac{e^{x}-e^{-x}}{2} \\ 
		\cosh(x) &= \frac{e^{x}+e^{-x}}{2} \\
		\tanh(x) &= \frac{\sinh(x)}{\cosh(x)}
	\end{split}
\end{equation*}

\begin{enumerate}[label=\alph*)]
	\item Erläutern Sie, wieso der Sinus-Hyperbolikus eine ungerade und der Kosinus-Hyperbolikus eine gerade Symmetrie aufweisen.

	\item Vergleichen Sie die Symmetrie mit der Sinus- und der Kosinus-Funktion.
	
	\item Plotten Sie alle 3 Hyperbolikus-Funktionen (z.B. mit GeoGebra oder Python).
	
\end{enumerate}

\end{ExerciseBlock}
\UseExerciseSolution{functions_hyperbolic_functions}
